\section{Wilson-loop hierarchies and the interacting experiment}\label{app:hierarchy}
\subsection{Continuous storage, memory, and projection error}
This appendix records the physical hierarchy from which the
source investigation developed. Its continuous statements
assume a positive regulated measure on smooth Hermitian
connections and bounded multiplication generators with
integrable operator norm on the capacity interval. They are
not asserted for unrenormalized continuum loops.

Let $q(t)\in\R^2$ be the tip of a planar trace starting at
the origin. Close the trace by its straight return chord.
If $U_t$ is transport along the trace and $R_t(r)$ along
the chord from 0 to $rq(t)$, define
\begin{equation}
Q_t=R_t(1)^{-1}U_t,\quad
J_t=\int_0^1r\,\widehat F_{12,t}(r)\dd r,\quad
k_t=q_1\dot q_2-q_2\dot q_1 .
\end{equation}
The hat denotes curvature transported back along the chord.
Variation of parallel transport, or differentiation of its
path-ordered equation with the endpoint terms cancelled,
gives $\dot Q_t=\ii k_tJ_tQ_t$. The moment $J_t$ is Hermitian
and traceless. Put $\mathsf A_tX=\ii k_tJ_tX$ in the
positive trace-normalized configuration Hilbert space,
$e=I$, $\Psi_t=Q_t$, and $W(t)=\ip e{\Psi_t}$.
The generator is skew-adjoint and the propagator is unitary.
For $P=|e\rangle\langle e|$ and $z=(I-P)\Psi$,
\begin{equation}
|W|^2+\norm z^2=1,\qquad
\dot W=-\ip{v_t}{z_t},\quad
\dot z_t=v_tW+(I-P)\mathsf A_t(I-P)z_t,\quad v_t=\mathsf A_te .
\end{equation}
The first identity is orthogonal decomposition. It describes
storage, not monotone dissipation: a two-dimensional unitary
rotation returns all amplitude. It also concerns one
prepared state, not arbitrary-input passivity.

The compressed bounded skew-adjoint generator has a unitary
propagator $\mathsf R(t,s)$. Eliminating z with $z_0=0$
gives the exact memory equation
\begin{equation}
\dot W(t)=-\int_0^t\Gamma(t,s)W(s)\dd s,\qquad
\Gamma(t,s)=\ip{v_t}{\mathsf R(t,s)v_s}.
\end{equation}
If $w_t=\mathsf R(t,0)^*v_t$, then
$\Gamma(t,s)=\ip{w_t}{w_s}$ for both time orders, so the
full two-time kernel is positive definite and
$\norm{z_T}^2=\norm{\int_0^Tw_sW(s)\dd s}^2$.
Writing $g(t)=\mathbb E\tr(J_t^2)/N_c$ also gives
$|\Gamma(t,s)|\le|k_tk_s|\sqrt{g(t)g(s)}$.
The causally truncated kernel is not thereby a positive
self-adjoint operator, and no stationarity in capacity
has been proved. This is an explicit time-dependent
projection identity in the spirit of \cite{mori}; evaluating
the complementary propagator can remain as difficult as
the original dynamics.

For a differentiable linearly independent family
$B(t)c=\sum_jc_jb_j(t)$ containing e, let
$G=B^*B$, $M=B^*\mathsf AB$, $C=B^*\dot B$.
Galerkin orthogonality gives
\begin{equation}
G\dot c=(M-C)c,\qquad
\frac{\dd}{\dd t}(c^*Gc)=0 .
\end{equation}
Indeed $M^*=-M$ and $\dot G=C+C^*$; omitting the moving
connection C invalidates conservation. For
$D=\dot B-\mathsf AB$ and $\Pi=BG^{-1}B^*$, the actual
residual is
\begin{equation}
r=(I-\Pi)Dc,\qquad
\norm r^2=c^*(D^*D-D^*BG^{-1}B^*D)c .
\end{equation}
This is a nonnegative Gram Schur complement. Singular
families must first be quotiented by their null space.
If the initial vector is represented exactly, Duhamel's
formula and unitarity imply
\begin{align}
\norm{\Psi_t-B(t)c(t)}&\le\int_0^t\norm{r_s}\dd s,\\
|W(t)-\ip e{B(t)c(t)}|
&\le\int_0^t\left(\int_s^t|k_v|\sqrt{g(v)}\dd v\right)
                                          \norm{r_s}\dd s.
\end{align}
For the second inequality use $\ip e{r_s}=0$ and subtract
e from the backward-propagated readout. Its difference
has norm at most the inner integral. Initial error adds
its norm to both bounds.

For the two-observable family $e,J/\sqrt g$, assume $g>0$
and put $\mu_j=\mathbb E\tr(J^j)/N_c$ and
$d=\mathbb E\tr(\dot J^2)/N_c$. The internal moving
connection vanishes and
\begin{equation}
\frac{\dd}{\dd t}\binom{x}{y}
=\ii k\begin{pmatrix}0&\sqrt g\\\sqrt g&\mu_3/g\end{pmatrix}
\binom{x}{y},\qquad (x(0),y(0))=(1,0).
\end{equation}
The orthogonal derivative and curvature residuals are
$R_D=\dot J-(\dot g/(2g))J$ and
$R_C=J^2-ge-(\mu_3/g)J$. Their inner product is real,
so the cross term in $\norm{R_D-\ii kR_C}^2$ vanishes:
\begin{equation}\label{eq:two-mode-residual}
\norm{r_2}^2=\frac{|y|^2}{g}
\left[d-\frac{\dot g^2}{4g}
+k^2\left(\mu_4-g^2-\frac{\mu_3^2}{g}\right)\right].
\end{equation}
Both summands in brackets are nonnegative squared
orthogonal residuals. The displayed plus sign is
essential; it repairs a missing addition sign in the
preliminary note's displayed equation (5.8).
For SU(2), $\mu_3=0$, so
$x=\cos\vartheta$, $y=\ii\sin\vartheta$, with
$\vartheta=\int_0^tk_s\sqrt{g(s)}\dd s$.
Although $J^2$ is scalar configuration by configuration,
its coefficient fluctuates, and $\mu_4-g^2$ need not vanish.

For the linear driver $u(t)=at$, the regulated tip expansion
gives $k_t=-(2a/3)\sqrt t+O(t^{3/2})$ and
$J_t=F_{12}(0)/2+O(\sqrt t)$. If
$C_0=\mathbb E\tr(F_{12}(0)^2)/N_c>0$, then
\[
\vartheta=-\frac{2a\sqrt{C_0}}9t^{3/2}+O(t^2),\qquad
W_2=1-\frac{2a^2C_0}{81}t^3+O(t^{7/2}).
\]
Under $\norm{R_D}=O(t^{-1/2})$ and $\norm{R_C}=O(1)$,
the refined output bound is $O(t^{7/2})$, while the
whole-state bound is only $O(t^2)$.
One sufficient regulated set of moment hypotheses is
finite $\mathbb E K_1^2$ and $\mathbb E K_0^4$, where
$K_0$ bounds curvature and $K_1$ its directional covariant
derivative in the swept region. Chord transport and the
ordered commutator term give
\begin{equation}
g\le\tfrac14\mathbb EK_0^2,\quad
\mu_4\le\tfrac1{16}\mathbb EK_0^4,\quad
d\le\mathbb E\left(|\dot q|K_1/3+|k|K_0^2/4\right)^2.
\end{equation}
The factors follow from $\int_0^1r\dd r=1/2$,
$\int_0^1r^2\dd r=1/3$, and the ordered-commutator
bound using $\norm{J(r)}\le K_0r^2/2$.
No regulator-uniform version is established.

\subsection{Exact discrete contour steps}
For actual lattice loops $Q_n$, multiplication by
$\mathsf U_n=Q_{n+1}Q_n^*$ is unitary and maps $Q_n$ to
$Q_{n+1}$. It uses no fractional link powers.
Relative to e and its orthogonal complement write the
blocks as $a_n,b_n,c_n,d_n$. With initial complement zero,
\begin{equation}
W_{n+1}=a_nW_n+
\sum_{j<n}b_nd_{n-1}\cdots d_{j+1}c_jW_j .
\end{equation}
This is exact block elimination; resetting the complement
would change the evolution.

For moving finite families $B_n$ containing e, set
$G_n=B_n^*B_n$ and $T_n=B_{n+1}^*\mathsf U_nB_n$.
Orthogonal projection gives
\begin{equation}
c_{n+1}=G_{n+1}^{-1}T_nc_n,\quad
\rho_n^2=\norm{B_{n+1}c_{n+1}-\mathsf U_nB_nc_n}^2
=c_n^*(G_n-T_n^*G_{n+1}^{-1}T_n)c_n .
\end{equation}
Here retained squared norm decreases by exactly $\rho_n^2$;
this discrete compression loss differs from the continuous
Galerkin conservation above. Telescoping unitary propagation
gives whole-state error at most $\sum_{j<N}\rho_j$.
Writing $K_{Nj}=\ip{Q_N}{Q_j}$, a sharper scalar bound is
\begin{equation}
|W_N-\widehat W_N|\le\sum_{j<N}\eta_{Nj}\rho_j,\qquad
\eta_{Nj}^2=2-2\Ree K_{N,j+1},
\end{equation}
where residual j lives at step $j+1$.
The factor is at most
$\min(2,\sum_{k=j+1}^{N-1}\sqrt{2-2\Ree K_{k+1,k}})$.
These follow by subtracting e from the exact backward
readout at that step, which is orthogonal to the residual.

A computable backward family gives another bound. Start
with $\chi_N=e=B_Nh_N$ and recurse
$h_j=G_j^{-1}T_j^*h_{j+1}$. Its backward residual has norm
\begin{equation}
\delta_j^2=h_{j+1}^*(G_{j+1}-T_jG_j^{-1}T_j^*)h_{j+1}.
\end{equation}
Backward telescoping bounds the error in the readout at
step $j+1$ by $\sum_{k=j+1}^{N-1}\delta_k$.
Since the forward residual is orthogonal to that retained
readout, one obtains
\begin{equation}\label{eq:fb-bound}
|W_N-\widehat W_N|
\le\sum_{j<N}\rho_j\sum_{k=j+1}^{N-1}\delta_k .
\end{equation}
All are inequalities for true ensemble Gram entries.
Floating Monte Carlo evaluation does not turn them into
certified ensemble bounds.

\subsection{The first interacting finite-slab test}
The completed experiment used the fixed SU(2) $6^3$ by
5 open-time Wilson lattice at $\beta=1.6$. Geometry was
specified before sampling: linear driver $u(t)=t$,
capacities $0,1,2.8,6$, nearest-vertex trace prefixes and
straight return chords. The signed lattice areas were
$0,-1,-2,-6$. The continuum tip obeys
$\dot q=1-2/q$ and
$q=2\ii\sqrt t+(2/3)t-(\ii/18)t^{3/2}+t^2/135+\cdots$.
Integration in $\sqrt t$ used maximum step $10^{-4}$,
started at $10^{-5}$. Halving the mesh left all words
unchanged; the independent equation
$q+2\log(1-q/2)=t$ had residual below $3.8\,10^{-12}$,
with tip refinement differences below $3.4\,10^{-12}$.

Four seeded chains, two cold and two Haar starts, used
600 burn-in sweeps and 2400 production sweeps each,
recording every sixth sweep: 1600 total configurations.
Each sweep used three proposals per link, with uniformly
distributed Pauli axis and angle in $[-1.2,1.2]$;
acceptance was approximately 0.599. The 216 spatially
translated loops were averaged within each configuration,
not counted as independent samples. Correlated batch
bootstrap used 400 resamples and block lengths 5, 10,
20, and 40; the principal quoted choice was 20.
Estimated autocorrelation scales of 0.50--0.76 measurement
intervals do not prove mixing.

The measured plaquette was $0.385525\pm0.000207$.
The loop expectations were
\begin{center}
\begin{tabular}{ccc}
\toprule
Capacity & Wilson mean & Estimated standard error\\
\midrule
1 & 0.388006 & 0.000716\\
2.8 & 0.150381 & 0.000939\\
6 & 0.002404 & 0.000891\\
\bottomrule
\end{tabular}
\end{center}
The tested local family used an anti-Hermitian transported
plaquette-curvature average B with $\norm B_{\rm op}\le1/2$.
The retained families were $I$; $I,B$; and $I,B,v^2I$,
where $B^2=-v^2I$ in SU(2). The unknown completed loop
was not included as a basis vector.
At capacity 2.8 the results were
\begin{center}
\begin{tabular}{rrrr}
\toprule
Dimension & Predicted mean & Difference from direct & Output bound\\
\midrule
1 & 0.150548 & $+0.000168\pm0.000746$ & 0.849452\\
2 & 0.135442 & $-0.014939\pm0.000666$ & 0.800524\\
3 & 0.136709 & $-0.013671\pm0.000678$ & 0.798915\\
\bottomrule
\end{tabular}
\end{center}
The first-step residual norms were approximately
0.92166, 0.86983, 0.86949; retained squared norms were
0.15055, 0.24340, 0.24398. Most state norm was lost.
The near scalar coincidence for dimension one is not
evidence of factorization, and larger nested families
need not improve a later scalar readout. These families
did not establish useful closure.

Earlier prescribed-background controls were more favorable:
for a noncommuting three-background SU(2) ensemble with
weights $0.2,0.3,0.5$, slope 0.8 and capacity 0.4225,
the direct loop was 0.9996561469, the two-mode result
0.9996420564, and the refined output bound about
$8.00\,10^{-5}$ versus observed error $1.41\,10^{-5}$.
The whole-state bound was about 0.00707.
This demonstrates the possible value of the projection
bound in a controlled ensemble, not its efficacy in YM.
The ensemble was not reflection/parity invariant.

The sampler sources, geometry rules, seeds, checks, and
small record are retained in \filepath{numerics/}.
Raw chains and executables are regenerable scratch
outputs, not a permanent deposited data archive.
No uncertainty estimate quoted here is an interval
certificate, a continuum extrapolation, or evidence
for an arithmetic source identity.
